\subsection{问题描述：波动电价场景下的购电策略}
在问题三日内合同修正调度框架基础上，问题四将固定购电电价替换为分时波动市场电价。0时刻无法预知全天真实电价，仅能依托历史时序数据构建电价预测；电价预测偏差会造成日前购电合同与储能调度方案偏离最优工况，带来应急购电、闲置合同等额外经济损失。
本节对比两种合同模式：Q4-2采用固定日前合同，0时刻签约后全天不可修改；Q4-3支持日内合同调整，在6:00、12:00、18:00基于更新信息修正剩余时段购电计划。两种方案均采用同一套近30日中位数电价预测基线，保证对比口径统一，用于量化日内合同调整机制在电价不确定性场景下的经济效益。所有决策严格遵守非前视约束：仅能使用决策时刻之前可观测数据，未来真实电价、源荷信息不参与优化求解，仅用于事后费用结算。

\subsection{基于历史信息的电价预测}
采用四种历史基准方法构建0时刻电价预测基线：昨日同时段、近7日均值、近30日均值、近30日中位数。对每个目标时段，选取对应时段历史电价进行统计估计。表\ref{tab:q4_price_pred_err}汇总四种预测方案的MAE、RMSE、尖峰MAE与最大绝对误差。

\begin{table}[htbp]
\centering
\caption{四种0:00基准价格预测方法误差}
\label{tab:q4_price_pred_err}
\begin{tabular}{lcccc}
\toprule
基准方法 & MAE & RMSE & 尖峰MAE & 最大绝对误差 \\
\midrule
昨日同时段 & 0.08433 & 0.12384 & 0.09226 & 0.53200 \\
近7日均值 & 0.08267 & 0.10527 & 0.09686 & 0.49669 \\
近30日均值 & 0.08748 & 0.11351 & 0.10835 & 0.54015 \\
近30日中位数 & 0.08328 & 0.12074 & 0.09880 & 0.62360 \\
\bottomrule
\end{tabular}
\footnotesize{单位：元/kWh。}
\end{table}

本文选取**近30日中位数**作为统一预测基线，该选择不是样本内最优预测器，而是为Q4-2、Q4-3提供完全一致的价格输入，隔离合同调整机制本身带来的收益。

设0时刻的近30日中位数价格基线为$\hat{c}_{d,t}^{(0)}$；日内$k$时刻之前已完成时段的价格误差为$e_{d,j}^c = c_{d,j}-\hat{c}_{d,j}^{(0)}$。利用已观测历史误差加权平均，随目标时段与当前时刻距离指数衰减，构造日内修正后的电价预测：
\begin{equation}
\hat{c}_{d,t|k} = \max\left\{10^{-4},\hat{c}_{d,t}^{(0)} + \bar{e}_{d,k}^c \exp\left[-(t-k)/18\right]\right\},\quad t \ge k.
\label{eq:q4_price_update}
\end{equation}
式中下限$10^{-4}$用于防止数值求解出现非负价格。该公式仅使用$k$时刻前已经观测到的价格数据，未来真实电价不参与决策，仅在结算阶段使用。

四种基准预测方法的误差排序与年度总费用排序并不完全一致。MAE对所有时段等权重，但运行费用同时受负荷、应急购电、储能状态共同影响。为保证对比公平，Q4-2、Q4-3统一使用近30日中位数电价预测；沿用问题二的负荷预测与官方光伏预报方案，避免负荷预测方法变动干扰价格与合同机制的对比。

\begin{figure}[htbp]
\centering
\includegraphics[width=0.85\textwidth]{figures/q4-price-diagnostics.png}
\caption{波动电价预测的拟合与误差分布}
\label{fig:q4_price_pred}
\end{figure}

\subsection{调度模型与两种合同模式}
Q4-2与Q4-3的核心差异在于购电合同是否允许日内调整。
\begin{itemize}
    \item \textbf{Q4-2（固定合同模式）}：在0时刻一次性生成全天购电合同，合同一经确定全天保持不变；日内仅依靠储能充放电与应急购电平抑功率缺口，不支持合同修改。
    \item \textbf{Q4-3（可调合同模式）}：0时刻生成初始购电合同，在6:00、12:00、18:00三个时点，利用更新的电价、光伏与负荷信息，对尚未执行的时段重新优化购电计划，已执行时段不可回溯修改。
\end{itemize}
两种模式在滚动LP优化中，均使用预测电价、预测负荷与光伏安排调度；合同调整沿用问题三的结算规则，仅相对0时刻原始计划进行一次结算，最终运行费用使用附件4的真实电价计算。

\subsection{仿真结果与分析}
\subsubsection{固定合同与可调合同对比}
两种模式全年仿真汇总结果如表\ref{tab:q4_two_mode}。
\begin{table}[htbp]
\centering
\caption{问题四两种模式的年度结果}
\label{tab:q4_two_mode}
\begin{tabular}{lccc}
\toprule
模式 & 总费用/元 & 应急购电量/kWh & 更新次数 \\
\midrule
Q4-2 固定合同 & 17378687.69 & 570732.66 & 0 \\
Q4-3 可调合同 & 15465962.60 & 210394.02 & 1002 \\
\bottomrule
\end{tabular}
\end{table}

由表可见，Q4-3相比Q4-2年度总费用减少1912725.09元，降幅11.01\%；应急购电量下降360338.64 kWh，降幅63.14\%。结果表明，在分时电价波动场景下，日内合同调整的价值体现在两方面：一是直接削减功率缺口对应的高价应急购电，二是通过储能跨时段优化，合理安排各时段购电量，降低电价波动带来的成本。

\subsubsection{不同价格基准与理想信息下界分析}
为验证结论鲁棒性，在四种历史电价预测基线、以及0时刻预知全部真实电价的理想信息场景下，分别求解Q4-2与Q4-3，结果如表\ref{tab:q4_price_baseline_oracle}。
\begin{table}[htbp]
\centering
\caption{问题四价格基准方法与理想信息结果}
\label{tab:q4_price_baseline_oracle}
\begin{tabular}{lcc}
\toprule
价格基准方法 & Q4-2 费用/元 & Q4-3 费用/元 \\
\midrule
昨日同时段 & 17402307.33 & 15485837.98 \\
近7日均值 & 17390589.20 & 15454638.40 \\
近30日均值 & 17400942.40 & 15455190.24 \\
近30日中位数 & 17378687.69 & 15465962.60 \\
0:00已知实际电价 & 17300085.39 & 15383049.81 \\
\bottomrule
\end{tabular}
\end{table}

四种历史基准方法下，Q4-2费用极差为23619.64元，Q4-3费用极差仅31199.58元，说明**合同调整机制带来的收益并不依赖某一种价格预测器**。
将全知理想电价场景作为下界，因无法提前获取真实电价带来的额外费用损失：Q4-2损失78602.31元，Q4-3损失82912.80元，仅占各自总费用的0.45\%与0.54\%。这说明波动电价带来的成本变化，不能全部归因于电价预测误差，同时包含电价水平、尖峰位置、光伏与负荷缺口的耦合作用。

在四种基准方案下，Q4-3费用均低于Q4-2，节约幅度稳定在约190万元。对比不同预测方法之间的费用差异，合同更新机制带来的收益，远大于更换简单电价预测器带来的收益。

\begin{figure}[htbp]
\centering
\includegraphics[width=0.85\textwidth]{figures/q4-outcome.png}
\caption{合同调整权与价格信息对年度费用的影响}
\label{fig:q4_contract_price}
\end{figure}

\subsection{本章小结}
本章针对分时波动电价场景，构建基于历史时序的电价预测方法，搭建固定合同、日内可调合同两类微网购电调度方案，开展全周期仿真对比。得到结论如下：
\begin{enumerate}
    \item 在近30日中位数电价预测基线下，相比固定合同方案，支持日内调整的Q4-3年度总费用降低11.01\%，应急购电量下降63.14\%，日内合同调整可以显著对冲电价波动带来的经济损失；
    \item 在多种电价预测基准下，可调合同策略均稳定优于固定合同，合同调整机制的收益不依赖特定的电价预测模型；
    \item 与预知全部真实电价的理想场景对比，电价预测误差带来的额外损失占总成本比例很低，电价波动场景的运行成本，是电价时序、源荷缺口、储能调度、合同规则多因素耦合的结果。
\end{enumerate}

本章从电价不确定性角度验证了日内合同调整机制的价值，与问题三光伏预测不确定性的结论形成完整互补，共同论证了多信息更新下滚动合同修正策略在微网经济调度中的有效性。

